% ===================================================================== Related Work
\section{Related Work}
\label{sec:related}

\subsection{External knowledge in NLI}
\label{sec:rw-knowledge}

Lexical knowledge has long been used to decide whether a hypothesis word is licensed by the premise.
KIM~\citep{chen2018neural} feeds WordNet relations (synonymy, antonymy, hypernymy, hyponymy,
co-hyponymy) as pairwise features into the co-attention, local inference and composition steps of ESIM,
and \citet{glockner2018breaking} showed that models trained on SNLI fail on pairs that differ by a
single lexical substitution, precisely the cases where such knowledge should help.

With pretrained
Transformers, knowledge is usually injected either as text or as entity embeddings: K-BERT
\citep{liu2020kbert} appends knowledge-graph triples to the input and restricts their visibility in
self-attention, while KnowBert~\citep{peters2019knowledge} and ERNIE~\citep{zhang2019ernie} link
mentions to knowledge-base entity embeddings.

Graph-based approaches retrieve a subgraph for the pair
and encode it with a graph neural network whose output is fused with the sentence representation
\citep{kapanipathi2020infusing}. \todo{add the references of the KGAnet and KGNLI baselines.}

A common limitation is that the retrieved evidence is used as an additional feature vector: it does not
change \emph{how} the decision is reached, and its effect on strong pretrained encoders is often small.
In \model{}, typed relations bias the premise--hypothesis alignment and enter the messages that each
hypothesis word receives, and the resulting word-level evidence is combined by explicit logical rules.

\subsection{Structure inside pretrained encoders}
\label{sec:rw-structure}

A second line of work integrates graphs with pretrained language models more tightly.
QA-GNN~\citep{yasunaga2021qagnn} reasons jointly over a question--answer context node and a
knowledge-graph subgraph, GreaseLM~\citep{zhang2022greaselm} exchanges information between the
language model and a graph neural network at several layers, and structural adapters
\citep{ribeiro2021structadapt} insert graph convolutions between frozen Transformer blocks.

These methods
target question answering or generation over knowledge or semantic graphs. \model{} instead builds its
graph over the contextual states of the encoder itself, connecting premise and hypothesis words through
lexical relations, and keeps the encoder untouched: the graph operates as a separate expert whose output
is added to the classifier logits.

\subsection{Alignment and natural logic}
\label{sec:rw-alignment}

Alignment between premise and hypothesis is a classic ingredient of NLI, from phrase-based aligners
\citep{maccartney2008phrase} to the soft attention of the decomposable attention model
\citep{parikh2016decomposable} and ESIM~\citep{chen2017enhanced}, whose matching features
($[\vect{x};\vect{m};\vect{x}\odot\vect{m};|\vect{x}-\vect{m}|]$) we reuse in the update of
\cref{eq:update}.

Natural logic~\citep{maccartney2009extended} goes further and composes lexical
relations along an alignment, taking into account the monotonicity of the context; NaturalLI
\citep{angeli2014naturalli} uses it for large-scale inference, and recent work makes the composition
differentiable \citep{feng2020exploring} or trains it with reinforcement learning and lexical resources
\citep{feng2022neuro}.

These systems depend on external aligners or polarity parsers that are mostly
available for English. \model{} uses a learned soft alignment with a \emph{null} target and does not model
monotonicity explicitly; the weakness this causes on downward-monotone contexts is discussed in
\cref{sec:results}.

\subsection{Atomic and probabilistic inference}
\label{sec:rw-atomic}

The aggregation in \model{} is closest to models that decide NLI from the parts of the hypothesis.
\citet{glickman2005probabilistic} defined the probability that a text entails a hypothesis as the product,
over hypothesis words, of the probability that each word is entailed by some text word, an idea later
extended by \citet{shnarch2011probabilistic}.

In neural NLI, SLR-NLI~\citep{stacey2022logical} treats
hypothesis spans as logical atoms, predicts a label for each span from sentence-level supervision, and
derives the pair label with rules such as ``a contradiction requires at least one contradicting span'';
it retains in-distribution accuracy and improves robustness when training data is scarce.

EPR~\citep{wu2023weakly} aligns phrases of the two sentences, predicts a label per aligned pair and
induces the sentence label with differentiable fuzzy-logic formulas, and FGLR~\citep{stacey2024atomic}
applies the same principle to facts generated by a large language model.

\model{} shares the decision rule of these models (any conflict yields contradiction; full support
yields entailment) but differs in three respects. First, its atoms are contextual token nodes rather than
spans, phrases or generated facts, so no chunker or generator is needed. Second, the alignment and the
messages of these atoms are shaped by typed lexical relations.

Third, the structured distribution is fused with, rather than substituted for, the sentence-level
classifier, which keeps phenomena that the rules cannot express within reach of the encoder.
\todo{confirm in the SLR-NLI and EPR papers whether either combines its prediction with a sentence-level
classifier or uses external knowledge.}

\subsection{Robustness and diagnostic evaluation}
\label{sec:rw-robustness}

NLI models exploit annotation artifacts \citep{gururangan2018annotation} and shallow heuristics such as
lexical overlap \citep{mccoy2019right}, so in-distribution accuracy overstates their inferential ability.

Diagnostic sets address this by annotating pairs with the phenomenon they test: the GLUE/SuperGLUE
diagnostic~\citep{wang2018glue,wang2019superglue}, which we use as AX-b, covers lexical semantics,
predicate--argument structure, logic and world knowledge, and MED~\citep{yanaka2019can} shows that models
fine-tuned on standard data handle upward but not downward monotonicity. Following this line, we report
results on out-of-distribution and diagnostic test sets in addition to in-distribution accuracy, and
analyse AX-b by phenomenon.

\subsection{NLI for Portuguese}
\label{sec:rw-portuguese}

Portuguese NLI has been driven by shared tasks and translated resources: ASSIN \citep{fonseca2016assin}
annotates entailment and paraphrase in news, SICK-BR~\citep{real2018sickbr} translates SICK, and ASSIN~2
\citep{real2020assin2} builds on SICK-BR with a binary entailment label. InferBR \citep{inferbr2024} provides
three-way labels and is the main benchmark of this work.

Most recent systems fine-tune Portuguese or
multilingual encoders such as BERTimbau~\citep{souza2020bertimbau}, Albertina~\citep{rodrigues2023albertina}
and XLM-R~\citep{conneau2020unsupervised}, and cross-lingual evaluation typically relies on XNLI
\citep{conneau2018xnli}, which does not include Portuguese.

Lexical-semantic resources for Portuguese
are comparatively rich, from Onto.PT~\citep{goncalooliveira2014ontopt} to PT-LKB
\citep{goncalooliveira2018ptlkb}, which merges ten resources, but to our knowledge they have not been
integrated into the inference mechanism of Transformer-based NLI models. \todo{check this claim with a
targeted search before submission.} \model{} uses PT-LKB in this way, and we also test whether the effect
transfers between English and Portuguese.
